\section{基本情况} \anon{杨林森}，\anon{男}，\anon{河北张家口人}，\anon[XXX ]{1998年1月出生}，\anon[XXX ]{西安电子科技大学}\anon[XXX ]{广州研究院}\anon{集成电路工程}专业2022级硕士研究生。 
\section{教育背景}
\begin{edubg}
\anon{2017.09～2021.07} & \anon{华北理工大学}，本科，专业:\anon{电子科学与技术}\\
\anon{2022.09～ }& \anon{西安电子科技大学}，硕士研究生，专业:\anon{集成电路工程}\\
\end{edubg}

\section{攻读硕士学位期间的研究成果}

\subsection{申请（授权）专利}
\begin{resresult}
	\item \anon[（第一学生发明人）]{袁嵩, \textbf{\anon{杨林森}}, 刘启帆等}.一种用于降低开关损耗的精确谷底开关电路:中国, 20241\\0228018.0
	\item \anon[（第一学生发明人）]{袁嵩, \textbf{\anon{杨林森}}, 刘启帆等}.一种用于提高电源能量传递效率的退磁时间逐步逼近电路:中国, 202510183610.8
\end{resresult}

%\begin{resresult}
%	\item 第一学生发明人.一种用于降低开关损耗的精确谷底开关电路:中国,2024102\\28018.0
%	\item 第一学生发明人.一种用于提高电源能量传递效率的退磁时间逐步逼近电路:中国,%202510183610.8
%\end{resresult}
%\subsection{学术论文}
%\item \anon[（第二学生作者）] {Qifan Liu, Xi Jiang, Xiaowu Gong, Song Yuan, \anon{Linsen Yang},Ying Wang}.A Novel Control Scheme for Light Load Power Reduction of Flyback Converter[J]. IEEE Transactions on Power Electronics, 2024. doi:10.1109/TPEL.2024.3489633
%\item Liu Q, Jiang X, Yang L, et al. Adaptive Soft Switching Asymmetrical Half Bridge Flyback Converter With Wider Voltage Range[J]. IEEE Journal of Emerging and Selected Topics in Power Electronics, 2025. doi:10.1109/JESTPE.2025.3540137
\subsection{参与科研项目及获奖}
\begin{resresult}
	\item 国家自然科学基金, 多域物联网系统数据安全关键技术研究, 2023.01-2027.12, 在研, 参与。
	\item 竞赛项目，第七届中国研究生创“芯”大赛，高精度RAMP设计，团队优秀奖，2024年8月。
	\item 2024年获西安电子科技大学硕士研究生一等学业奖学金。
\end{resresult}
